% ICLR Paper Skeleton
% Replace placeholders marked with <...> before use
\documentclass{article}

% ICLR style — update year as needed
\usepackage{iclr2025_conference}

\usepackage[utf8]{inputenc}
\usepackage[T1]{fontenc}
\usepackage{hyperref}
\usepackage{url}
\usepackage{booktabs}
\usepackage{amsfonts}
\usepackage{amsmath}
\usepackage{amssymb}
\usepackage{nicefrac}
\usepackage{microtype}
\usepackage{graphicx}
\usepackage{xcolor}

\newcommand{\ours}{\textsc{<MethodName>}}
\newcommand{\todo}[1]{\textcolor{red}{[TODO: #1]}}

\title{<Paper Title>}

% Authors — anonymized for submission
\author{Anonymous}

\begin{document}

\maketitle

\begin{abstract}
% 4-sentence structure:
% S1: Problem context and importance
% S2: Key limitation of existing approaches
% S3: Our approach and key insight
% S4: Main results and significance
\todo{Write abstract — 150-250 words}
\end{abstract}

%% ============================================================
\section{Introduction}
\label{sec:intro}

\todo{Broad motivation}

\todo{Specific problem statement}

\todo{Prior work limitations}

\todo{Our approach — key insight}

Our contributions are summarized as follows:
\begin{itemize}
    \item \todo{Contribution 1}
    \item \todo{Contribution 2}
    \item \todo{Contribution 3}
\end{itemize}

%% ============================================================
\section{Related Work}
\label{sec:related}

\paragraph{<Topic 1>.}
\todo{Related work group 1}

\paragraph{<Topic 2>.}
\todo{Related work group 2}

\paragraph{<Topic 3>.}
\todo{Related work group 3}

%% ============================================================
\section{Preliminaries}
\label{sec:prelim}

% ICLR papers often include a Preliminaries section
\todo{Background, notation, problem setup}

%% ============================================================
\section{Method}
\label{sec:method}

\todo{Method overview — reference Figure 1}

\subsection{<Component 1>}
\label{sec:component1}
\todo{First key component with formal description}

\subsection{<Component 2>}
\label{sec:component2}
\todo{Second key component}

\subsection{Theoretical Analysis}
\label{sec:theory}
% ICLR values theoretical grounding
\todo{Theoretical justification or analysis}

%% ============================================================
\section{Experiments}
\label{sec:experiments}

\subsection{Setup}
\label{sec:setup}
\todo{Datasets, baselines, metrics, implementation details}

\subsection{Main Results}
\label{sec:main_results}
\todo{Main comparison and analysis}

\subsection{Ablation Study}
\label{sec:ablation}
\todo{Component-wise analysis}

\subsection{Analysis}
\label{sec:analysis}
\todo{Qualitative results, visualizations, scaling behavior}

%% ============================================================
\section{Discussion and Conclusion}
\label{sec:conclusion}

\todo{Summary, limitations, broader impact, future work}

% Reproducibility Statement (encouraged by ICLR)
\subsubsection*{Reproducibility Statement}
\todo{Describe steps taken to ensure reproducibility}

% Ethics Statement (if applicable)
% \subsubsection*{Ethics Statement}
% \todo{Ethical considerations}

\bibliography{references}
\bibliographystyle{iclr2025_conference}

%% ============================================================
\appendix
\section{Proof Details}
\label{app:proofs}
\todo{Full proofs}

\section{Additional Experiments}
\label{app:experiments}
\todo{Additional results, hyperparameter sensitivity}

\section{Implementation Details}
\label{app:implementation}
\todo{Full implementation and training details}

\end{document}
